\frametitle{The Correction Is the Complement Projector}
\footnotesize

\begin{block}{Proposition}
Let $A_c=P^{\mathsf T}AP$ be the Galerkin coarse operator and let $R=P^{\mathsf T}$. Then
\[
PA_c^{-1}RA \;=\; P\bigl(P^{\mathsf T}AP\bigr)^{-1}P^{\mathsf T}A \;=\; \PiC ,
\qquad\text{and hence}\qquad
C_h \;=\; I-\PiC \;=\; \PiF ,
\]
where $\PiC$ is the $A$-orthogonal projection onto $\operatorname{range}(P)$ and
$\PiF=I-\PiC$ its $A$-orthogonal complement.
\end{block}

{\scriptsize\emph{Proof.} The middle factor is precisely the stated expression.
$\PiC^2=\PiC$ follows by cancelling $(P^{\mathsf T}AP)^{-1}(P^{\mathsf T}AP)$, and
$\PiC^{\mathsf T}A=A\PiC$ is immediate from the symmetry of $A$ and $A_c$.
\hfill$\blacksquare$}

\vspace{1.5mm}
Two things are worth stating plainly.
\begin{itemize}
\item It is an \alert{identity, not an approximation}: it holds exactly and is checkable in
the same code that produced the experiment, with no new data.
\item It is the reason the two-level analysis below is possible at all --- because
$C_h=\PiF$ exactly, the two-grid error propagator is
\[
\ETG=S_{\mathrm{post}}\bigl(I-PA_c^{-1}RA\bigr)S_{\mathrm{pre}}
=S_{\mathrm{post}}\,\PiF\,S_{\mathrm{pre}},
\]
with no error term to account for. Anything in the cycle that is not a smoothing sweep is a
projection, and a projection is something one can reason about exactly.
\end{itemize}

\bre
The identity uses an \emph{exact} solve of the coarse problem. In practice the coarsest
level is small enough to be solved directly, so this is the configuration the solver
actually runs. Where the recursion continues to further levels, the two-grid factor becomes
the appropriate bound rather than an identity --- the classical distinction between a
two-grid analysis and a V-cycle analysis \cite{Trottenberg}, and exactly the gap Stage~III
has to close.
\ere
